\subsection{Positive initial accelerations}

\begin{lemma}[Finite initialization and layerwise activity]
\label{fl:initial-activity}
Under the additional theorem's hypotheses, each layer's initialized forward
fields, first backward derivatives (also before activation gates), and
preactivation and feature accelerations have deterministic joint empirical
row-law limits with convergent second moments. This includes the full
first-layer weight and acceleration rows. Every hidden block's squared
acceleration in inverse-mobility norm and every layer's squared feature
acceleration in neuron RMS converge to strictly positive constants.
The corresponding squared norms are uniformly integrable.
\end{lemma}

\begin{proof}
All notation below is local and all unmarked network fields are at zero.
Put $k=2/m$,
$\|q\|_n=\|q\|_F/\sqrt n$ for any array with $n$ rows, and define
\begin{align}
 s&=\sum_a y_a h_a^{(L)},&
 B_a^{(L)}&=s\odot\phi_L'(z_a^{(L)}),\nonumber\\
 B_a^{(\ell)}&=\phi_\ell'(z_a^{(\ell)})\odot
                  W_0^{(\ell+1)\top}B_a^{(\ell+1)},
 &V_\ell&=\ddot W^{(\ell)}(0),\nonumber\\
 R_a^{(\ell)}&=\ddot z_a^{(\ell)}(0),&
 A_a^{(\ell)}&=\ddot h_a^{(\ell)}(0).
 \label{fl:initial-fields}
\end{align}
Here $W_0^{(\ell)}=W^{(\ell)}(0)$, and
$\dot\delta_a^{(\ell)}(0)=kB_a^{(\ell)}$.
Each layer's joint empirical row law of its forward fields, $B,R,A$, and
the pre-gated fields $s$ or $W_0^{(\ell+1)\top}B_a^{(\ell+1)}$
converges in probability in quadratic Wasserstein distance to a deterministic
law. At layer one this includes the whole root row
$g_i=W_{0,i:}^{(1)}\in\R^d$ and the acceleration row $V_{1,i:}$.
In particular, all same-layer normalized inner products converge.

For the mobility norms
$\|V_1\|_{\rm mob}=\|V_1\|_n$ and
$\|V_\ell\|_{\rm mob}=\|V_\ell\|_F$ for $\ell\ge2$, there are
deterministic $e_\ell,c_\ell>0$ such that
\begin{equation}
 \|V_\ell\|_{\rm mob}^2\longrightarrow e_\ell,
 \qquad \sum_a\|A_a^{(\ell)}\|_n^2\longrightarrow c_\ell
 \quad\hbox{in probability}.
 \label{fl:positive-accelerations}
\end{equation}
We also establish uniform integrability of these squared norms and uniform
integrability in probability of their squared coordinate tails.
Since all hidden velocities initially vanish, differentiating the physical
updates and the forward recursion gives exactly
\begin{align}
 V_1&=k^2\sum_a y_aB_a^{(1)}v_a^\top,&
 V_\ell&=\frac{k^2}{n}\sum_a y_aB_a^{(\ell)}
                    h_a^{(\ell-1)\top}\quad(\ell\ge2),\nonumber\\
 R_a^{(1)}&=V_1v_a,&
 R_a^{(\ell)}&=V_\ell h_a^{(\ell-1)}
                         +W_0^{(\ell)}A_a^{(\ell-1)},\nonumber\\
 A_a^{(\ell)}&=\phi_\ell'(z_a^{(\ell)})\odot R_a^{(\ell)}.
 \label{fl:acceleration-recursion}
\end{align}
We prove the needed row-law limits by a finite forward--backward--forward
Gaussian conditioning argument, retaining both directions of each matrix.

For one hidden matrix $W$, suppose its recorded actions are $WH=Z$ and
$W^\top U=T$, with fixed column counts. On invertible query Grams set
$P_H=H(H^\top H)^{-1}H^\top$ and likewise $P_U$. Conditional on the
recorded transcript, its remaining law is
\begin{equation}
 W=Z(H^\top H)^{-1}H^\top
    +U(U^\top U)^{-1}T^\top(I-P_H)
    +(I-P_U)\widetilde W(I-P_H),
 \label{fl:two-sided-conditioning}
\end{equation}
where $\widetilde W$ has independent $N(0,1/n)$ entries. Indeed
$U^\top Z=T^\top H$, so the displayed mean satisfies both constraints.
Their homogeneous solution space consists exactly of
$(I-P_U)E(I-P_H)$, and that mean is Frobenius-orthogonal to the space.
Independence of orthogonal Gaussian projections proves the formula.
Before a reverse call omit its absent query block.

The identity also holds for interlaced adaptive calls to the independent
matrices. Condition first on the first-layer roots. Given the answers so
far, the next query used here is fixed and reveals a linear projection of
one remaining Gaussian matrix subspace. Its orthogonal complement stays
Gaussian and independent of all other residual matrices. Induction on calls
proves the stated conditional law for the complete global transcript.

In particular, for the first reverse call write
$Q_n=H^\top H/n$, $C_n=Z^\top U/n$, $D_n=U^\top U/n$. It gives
\begin{equation}
 T\overset d=H Q_n^{-1}C_n+(I-P_H)\Xi D_n^{1/2},
 \qquad
 \E[\|P_H\Xi D_n^{1/2}\|_n^2\mid\mathrm{past}]
   =\frac{\operatorname{rank}H}{n}\operatorname{tr}D_n,
 \label{fl:reverse-conditioning}
\end{equation}
where $\Xi$ has independent standard Gaussian entries, independent of the
past. After this answer, a new forward query $A$ obeys
\begin{align}
 F_n&=H^\top A/n,&J_n&=T^\top A/n,&E_n&=A^\top A/n,\nonumber\\
 \Sigma_n&=E_n-F_n^\top Q_n^{-1}F_n,\nonumber\\
 WA&\overset d=ZQ_n^{-1}F_n
  +UD_n^{-1}(J_n-C_n^\top Q_n^{-1}F_n)
  +(I-P_U)\Xi\Sigma_n^{1/2}.
 \label{fl:acceleration-conditioning}
\end{align}
Now $\Xi$ is fresh after the reverse answer. Since
$\Sigma_n=A^\top(I-P_H)A/n\succeq0$, it may be singular and is never
inverted. The removed projection's conditional squared RMS is
$\operatorname{rank}(U)\operatorname{tr}(\Sigma_n)/n$.

Here are the complete limit rules for these finitely many calls. Appending
independent standard Gaussian rows to an array with a deterministic empirical
$\mathcal W_2$ limit appends independent Gaussian variables to that limit.
For bounded continuous tests, conditional variances of empirical averages
are $O(1/n)$ and their conditional means converge by the old empirical law.
The appended second moment converges by the Gaussian law of large numbers.
Weak convergence and convergence of second moments give $\mathcal W_2$
convergence. Same-layer empirical products therefore converge, since they
are continuous of at most quadratic growth. Gram inverses converge whenever
their limits are positive definite. Positive square roots converge even at
singular covariance limits: bounded subsequences have limits whose squares
equal the limiting covariance, and its positive square root is unique.
Thus multiplying the fresh Gaussian rows by these roots preserves joint
convergence. The projection errors above vanish in RMS in probability,
since their covariance traces stay bounded in probability. Coupling rows
by their indices shows that these errors do not change a $\mathcal W_2$ limit.

The coordinate maps are linear combinations, Lipschitz activations, and
$(z,q)\mapsto\phi_\ell'(z)q$. Each is continuous with at most linear
growth in its complete input tuple and therefore preserves $\mathcal W_2$
convergence: on an $L^2$ coupling, continuity gives convergence in probability,
and uniform integrability of squared inputs and the linear envelope give
$L^2$ convergence. Coefficients converging to constants are handled first
by subtracting their linear combinations and using bounded input RMS.

Start with the iid full rows $g_i\sim N(0,I_d)$. Their empirical laws
converge with second moments. Make all initial forward calls in order.
Conditional fresh Gaussian rows and the preceding rules yield the forward
tuple limits $Z^{(\ell)}\sim N(0,Q^{(\ell-1)})$ and
$H_a^{(\ell)}=\phi_\ell(Z_a^{(\ell)})$.
Lemma~\ref{fl:input-lemma} gives $Q^{(\ell)}\succ0$ for $\ell\ge1$.

Write $\mathsf B,\mathsf R,\mathsf A$ for limiting row variables of
$B,R,A$, and $D^{(\ell)}_{ab}=\E[\mathsf B_a^{(\ell)}
\mathsf B_b^{(\ell)}]$. The top response is the continuous, at-most-linear
map $\mathsf B_a^{(L)}=(\sum_b y_bH_b^{(L)})\phi_L'(Z_a^{(L)})$.
Its Gram is positive definite. If a linear combination with coefficients
$c_a$ vanishes, full Gaussian support and continuity imply
\[
 \Bigl(\sum_b y_b\phi_L(z_b)\Bigr)
 \Bigl(\sum_a c_a\phi_L'(z_a)\Bigr)=0
 \quad\hbox{for every }z\in\R^m.
\]
The first factor is not identically zero, since $y^\top Q^{(L)}y>0$.
The second vanishes on a nonempty open set, hence everywhere by analyticity.
Varying one coordinate and using the nonconstant derivative gives $c_a=0$.

Proceed downward. In \eqref{fl:reverse-conditioning} take
$H=[h_a^{(\ell)}]$, $Z=[z_a^{(\ell+1)}]$ and
$U=[B_a^{(\ell+1)}]$. All upper responses are already known from the
forward pass and earlier reverse calls; this is the first reverse call to
the present matrix. The preceding limit rules give its limiting answer
\[
 \sum_bH_b^{(\ell)}[(Q^{(\ell)})^{-1}C]_{ba}+\Gamma_a,
 \quad C_{ba}=\E[Z_b^{(\ell+1)}\mathsf B_a^{(\ell+1)}],
 \quad\Gamma\sim N(0,D^{(\ell+1)}),
\]
where $\Gamma$ is independent of the previous lower-layer tuple, including
$g$ at layer one. Multiplication by $\phi_\ell'(Z_a^{(\ell)})$ yields
$\mathsf B_a^{(\ell)}$, and conditional variance gives, for $c\ne0$,
\[
 c^\top D^{(\ell)}c\ge\lambda_{\min}(D^{(\ell+1)})
       \sum_a c_a^2\E[\phi_\ell'(Z_a^{(\ell)})^2]>0.
\]
Every marginal preactivation is nondegenerate Gaussian; a nonzero analytic
derivative has a discrete zero set. This proves the strict inequality and
completes both the backward row-law and positive-Gram inductions.

Construct the accelerations forward using \eqref{fl:acceleration-recursion}.
At layer one the joint limit retains the exact row formulas
\[
 Z_a^{(1)}=g^\top v_a,\quad
 V_{1,\mathrm{row}}=k^2\sum_a y_a\mathsf B_a^{(1)}v_a,\quad
 \mathsf R_a^{(1)}=V_{1,\mathrm{row}}^\top v_a,\quad
 \mathsf A_a^{(1)}=\phi_1'(Z_a^{(1)})\mathsf R_a^{(1)}.
\]
For later layers the learned-link contribution is exactly
$k^2\sum_b y_bB_b^{(\ell)}Q_{n,ba}^{(\ell-1)}$ and has its row-law
limit. Apply \eqref{fl:acceleration-conditioning} to the remaining term
with $A=[A_a^{(\ell-1)}]$ and the already recorded forward and reverse
queries. Its coefficients are same-layer second moments already obtained;
its only inverse Grams are the positive $Q^{(\ell-1)},D^{(\ell)}$.
This proves joint convergence of the new preactivation and feature
accelerations. There are exactly $3(L-1)$ hidden-matrix calls; each query
is determined by earlier answers. Inverses are used only on events of
probability tending to one, which suffices for these limits.

Uniform integrability follows without inverses. The variable
$M_n=1+\|W_0^{(1)}\|_n+\sum_{\ell\ge2}\|W_0^{(\ell)}\|_{\rm op}$
has bounded moments of every fixed order, uniformly in width. For the first
term use Jensen and Gaussian moments; for the others integrate the two-net
Gaussian tail used in Lemma~\ref{fl:bootstrap}. Linear activation growth,
bounded derivative gates and \eqref{fl:acceleration-recursion} bound every
array RMS and block mobility norm by a fixed polynomial in $M_n$.
Thus their squared norms have bounded moments above order one. This gives
uniform integrability of the energies, while the proved $\mathcal W_2$
convergence gives uniform integrability in probability of squared coordinate
tails. In particular all limits of squared norms are also $L^1$ limits.

Finally the exact transpose relation in \eqref{fl:initial-fields} gives
the partial-layer adjoint identity
\begin{equation}
 \sum_a\frac{y_a}{n}B_a^{(\ell)\top}R_a^{(\ell)}
       =\frac1{k^2}\sum_{j=1}^{\ell}\|V_j\|_{\rm mob}^2.
 \label{fl:partial-adjoint}
\end{equation}
Indeed the learned-link term is $\|V_\ell\|_{\rm mob}^2/k^2$ by
the formula for $V_\ell$, and the propagated term equals the left side
one layer lower; at the first layer only the learned-link term remains.
The squared mobility limits are
$e_\ell=k^4y^\top(Q^{(\ell-1)}\circ D^{(\ell)})y$.
For $Q=\sum_jq_jq_j^\top\succeq0$ and $D\succ0$,
$Q\circ D=\sum_j\operatorname{diag}(q_j)D\operatorname{diag}(q_j)
\succeq\lambda_{\min}(D)\operatorname{diag}(Q)$.
Every relevant $Q$ has positive diagonal, including $Q^{(0)}$, so
$e_\ell>0$. Passing to the established joint-law limits in
\eqref{fl:partial-adjoint} shows that some $\mathsf R_a^{(\ell)}$ is
nonzero in $L^2$. Since $\phi_\ell'(Z_a^{(\ell)})\ne0$ almost surely,
$\mathsf A_a^{(\ell)}=\phi_\ell'(Z_a^{(\ell)})\mathsf R_a^{(\ell)}$
cannot vanish for all $a$. Thus
$c_\ell=\sum_a\E|\mathsf A_a^{(\ell)}|^2>0$, as claimed.
\end{proof}
